\chapter{AHC：内稳态—情感调制层}

\noindent
\fbox{
\parbox{0.96\textwidth}{
\textbf{本章总体说明（理论定位）}\quad
在前述工作记忆有限容量理论、SPC 自我状态估计机制、TCE 认知动力学调节机制以及工作记忆运行相态理论建立之后，我们已经获得了一个最小的“观察者--控制器”认知闭环：
\[
SPC \rightarrow TCE \rightarrow CognitiveDynamics \rightarrow SPC .
\]
然而，这一闭环仍然隐含了一个过强假设：只要系统能够正确估计自身认知状态，控制器便可以直接根据这些状态产生合适的认知调节。真实的持续智能体并非如此。资源消耗、风险暴露、身体完整性、任务紧迫度、成功与失败、社会反馈、长期疲劳以及不确定性，并不会以完全相同的时间尺度和相同的语义形式直接作用于认知控制。它们首先形成一种具有积分、衰减、累积、恢复、饱和和相互抑制特性的中间动力状态，然后才改变注意、搜索、探索、收敛、记忆写入和行动准备。
因此，本章引入 AHC（Affective--Homeostatic Co-processor，内稳态--情感协处理器）。AHC 并不是为了在人工智能中复制人类情绪，也不构成关于机器主观意识的论断；它首先是一个\textbf{多源内部调制信号的中时间尺度动力学积分层}。其理论任务是把瞬时的身体、资源、风险、奖励、社会关系和认知压力转换为具有时间连续性的调制状态，使认知系统不仅能够回答“我现在发生了什么”，还能够形成“这种状态对我意味着什么，以及它应当以多大力度影响接下来的认知过程”的内部动力学。
本章只建立 AHC 本身的状态空间、动力学、输入输出关系和稳定性原则。SPC--AHC--TCE 的完整体验闭环将在后续章节进一步统一。
}
}

\section{为什么 SPC 与 TCE 之间仍然需要一个中间动力层}

在上一阶段的 AgentOS 原理论中，我们已经把 SPC 理解为一种内部状态观察器。它面对的不是单一变量，而是工作记忆负载、冲突程度、当前目标竞争、不确定性、循环程度、认知相态以及部分身体与资源状态。SPC 的任务可以抽象为
\[
\hat{x}^{self}_t
=
\mathcal O_{\mathrm{SPC}}
\left(
z_{\leq t}
\right),
\]
其中 \(z_{\leq t}\) 表示截至当前时刻能够被系统获得的内部与外部证据，\(\hat{x}^{self}_t\) 则是 Agent 对自身当前状态形成的有限后验估计。随后，TCE 根据这一估计产生控制量
\[
u_t
=
\pi_{\mathrm{TCE}}
\left(
\hat{x}^{self}_t,
G_t,
\Psi_t
\right),
\]
改变搜索温度、注意范围、探索强度、推理深度和认知节律。表面上看，这已经形成完整的 observer--controller 闭环。但如果我们把这一关系直接写成
\[
\hat{x}^{self}_t
\longrightarrow
u_t,
\]
就会遇到一个结构性问题：许多输入虽然可以在瞬间被观察，却不应该在瞬间等比例地改变认知控制。

例如一次短暂的计算资源下降，并不意味着系统立即进入长期保守模式；一次偶然失败，也不应该立刻改变整个 Agent 对某类任务的长期风险倾向；相反，如果能源持续下降、错误连续出现、风险信号不断积累，即使任何一个单独时刻的信号都不算极端，系统也应该逐渐提高警觉并收缩探索空间。换言之，认知控制所需要的不是一个纯粹的瞬时状态向量，而是一个经过时间积分以后形成的\textbf{状态意义场}。

设原始内部影响信号为
\[
r_t
=
[r_t^{1},r_t^{2},\ldots,r_t^{m}]^\top .
\]
如果直接控制，则有
\[
u_t = K r_t .
\]
这种结构对高频扰动高度敏感。更合理的方式是先建立一个中间状态 \(a_t\)：
\[
\tau_a \dot a_t
=
-a_t
+
F(r_t,a_t),
\]
再由
\[
u_t
=
\pi_{\mathrm{TCE}}
(a_t,\hat{x}^{self}_t,G_t)
\]
生成控制。这里的 \(a_t\) 不再是原始传感量，而是经过积分、衰减、交互和归一化之后形成的调制状态。AHC 的理论位置正是在这里产生。

因此，我们把最小认知控制链从
\[
SPC\rightarrow TCE
\]
扩展成
\[
\boxed{
State
\rightarrow
SPC
\rightarrow
AHC
\rightarrow
TCE
\rightarrow
CognitiveDynamics
}
\]
但应当立即指出，AHC 并不简单夹在两个模块之间形成机械串联。真实关系更接近一个耦合系统：
\[
\dot a
=
F_a(a,\hat{x}^{self},r),
\]
\[
u
=
F_u(a,\hat{x}^{self},G),
\]
\[
\dot x
=
F_x(x,u),
\]
其中 \(a\) 具有自己的时间尺度，并会持续影响未来多个时刻的控制决策。

这意味着 AHC 第一次在 AgentOS 中引入了一个重要概念：\textbf{当前状态的影响可以持续存在，而不要求原始刺激继续存在。} 一个短暂事件可以产生一个随时间逐渐衰减的内部动力轨迹；多个较弱事件则可能由于时间积累而形成显著影响。这种机制使持续智能体开始区别于简单的事件触发程序，也为后面研究认知恢复、压力积累、情感样连续性以及心理工程中的动力异常提供了可计算基础。

从更一般的智能动力学角度看，AHC 解决的是不同时间尺度之间的阻抗匹配问题。外部环境和身体可能产生毫秒级变化，工作记忆可能以秒到分钟变化，而认知策略不能对每个高频扰动都同步重构。AHC 因而承担一个低通、积分和非线性重编码作用：
\[
FastDisturbance
\rightarrow
MediumScaleInternalState
\rightarrow
CognitiveModulation .
\]
这一结构与生物内稳态、控制理论中的滤波状态以及神经调制系统存在可以比较的形式关系，但本书并不把这些生物机制直接复制到人工智能中。我们真正需要的是一个实现无关的动力学功能：\textbf{把大量瞬时、异质、局部的“对主体产生影响的信号”，转换成少量连续、稳定、能够参与认知调节的内部状态。}

因此，AHC 的引入并不是为了给 AgentOS 增加“情绪模块”，而是由多时间尺度控制本身强制推出的。如果没有这一层，SPC--TCE 很容易退化成瞬时、高增益、对噪声敏感的控制系统；有了这一层，Agent 才能够开始形成“持续状态”，并让过去一段时间中的风险、资源、成功、失败和社会关系以有记忆的方式共同影响当前认知。

\section{AHC 的状态空间：从离散标签到连续内稳态向量}

我们首先必须避免一个常见误区：把 AHC 设计成一组离散的“开心、悲伤、愤怒、害怕”标签。如果这样做，AHC 就会变成对人类情感词汇的表层模仿，而不是具有一般智能意义的动力结构。在 AgentOS 中，AHC 首先应该表示的是一组\textbf{连续、可测、具有动力学时间常数的功能状态}。

定义 AHC 状态为
\[
\boxed{
a(t)
=
[
a_1(t),
a_2(t),
\ldots,
a_n(t)
]^\top
\in
\mathcal A
}
\]
其中每一个维度并不要求与自然语言中的某种情绪一一对应，而是表示一种对认知和行为具有系统性调制作用的内部变量。根据此前讨论，可以构造一个基础状态集合：
\[
a(t)=
[
A_r,
T_h,
E_d,
F_g,
R_e,
S_a,
T_r,
U_s
]^\top ,
\]
分别表示 Arousal、Threat、Energy Drive、Fatigue、Reward Expectation、Social Affinity、Trust 与 Uncertainty Stress。这里的符号仅作为一种标准化起点，而不是宣称所有 Agent 必须具有完全相同的维度。

例如，Arousal \(A_r\) 可以表示系统整体激活准备程度。它并不是“兴奋”这一自然语言情绪，而更加接近：
\[
A_r
\sim
\text{readiness of cognitive and behavioral activation}.
\]
当 \(A_r\) 上升时，TCE 可以提高注意刷新频率、缩短响应周期或增加局部搜索资源；但如果 \(A_r\) 过高，也可能降低稳定推理能力。因此它通常不是单调收益变量。

Threat \(T_h\) 表示当前状态中潜在损失、不可逆错误或安全风险的聚合强度。它与外部“危险物体”的语义表示不同：前者描述这种风险对当前主体控制策略的影响程度。形式上可以写成
\[
T_h(t)
=
\Phi_T
\left(
p(\mathrm{failure}),
L(\mathrm{failure}),
C_{\mathrm{recover}})
\right),
\]
其中风险不仅与失败概率有关，也与损失和可恢复性有关。

Energy Drive \(E_d\) 与 Fatigue \(F_g\) 应被分别保留。前者表示资源继续投入的可用性与驱动力，后者则描述累计活动导致的运行负担。二者并不是简单互补：
\[
F_g \neq 1-E_d.
\]
例如系统可能同时处于高资源、高疲劳状态，也可能处于低资源但负载暂时较低的状态。因此不同内部变量必须允许部分独立。

Reward Expectation \(R_e\) 不是简单 reward 数值，而是对当前行动链未来正向结果可达性的内部估计：
\[
R_e(t)
=
\mathbb E
\left[
R_{t:t+H}
\mid
\mathcal I_t
\right].
\]
它影响系统是否继续投入资源、是否扩大探索以及是否保持当前策略。

Social Affinity 与 Trust 则用于处理社会智能环境中的两个不同问题。Affinity 更接近关系接近度、合作倾向和交互价值；Trust 更接近对另一主体预测可靠性和承诺可兑现程度的估计。一个 Agent 完全可能具有
\[
Affinity\uparrow,\qquad Trust\downarrow ,
\]
例如愿意帮助某个对象但并不依赖其信息。这个区分对于未来多 Agent 系统非常重要。

Uncertainty Stress \(U_s\) 则把纯粹的认识论不确定性与其产生的控制影响区分开来。SPC 可以估计
\[
H[p(x)]
\]
很高，但系统是否因此提高警觉、收缩行动空间，还取决于这一不确定性是否与当前目标和风险相关。因此：
\[
Uncertainty
\neq
UncertaintyStress.
\]
前者是认知状态，后者是这种认知状态经过目标、价值和风险解释之后形成的调制量。

由此我们得到一个重要原则：
\[
\boxed{
AHCState
\neq
SemanticDescription.
}
\]
AHC 状态首先是动力状态，而自然语言描述只是其某种高层可解释映射：
\[
Label_t
=
\mathcal L(a_t).
\]
例如一个状态向量可能被语言模型解释为“紧张但有动力”，但这种描述并不是底层变量本身。工程上必须坚持
\[
a_t \rightarrow Label_t
\]
而不是反过来用标签替代状态动力学。

进一步，为保证可计算性，我们通常把每个状态规范到有限区间：
\[
a_i\in[a_i^{\min},a_i^{\max}],
\]
例如
\[
a_i\in[0,1]
\]
或
\[
a_i\in[-1,1].
\]
但归一化只是一种表示方式，不代表所有维度具有相同物理意义。不同状态还需要拥有独立的时间常数：
\[
\tau_i,
\]
以及不同饱和阈值
\[
\theta_i^{sat}.
\]

因此，完整的 AHC 状态并不仅是一个值向量，而更接近：
\[
\boxed{
\mathcal A_t
=
\{
a(t),
\dot a(t),
\Sigma_a(t),
\tau,
\Theta_{sat}
\}
}
\]
其中 \(\Sigma_a\) 表示状态估计不确定性。这样，TCE 不仅知道“当前 Threat=0.7”，还可以知道 Threat 正在快速上升还是缓慢下降，以及这一估计到底有多可靠。

我们由此得到 AHC 的第一个标准工程定义：

\[
\boxed{
\text{AHC 是一个由有限维连续状态、独立时间常数、非线性相互作用和不确定性描述共同构成的中尺度内部调制动力系统。}
}
\]

这一表述有意避免把 AHC 拟人化。它使任何具有不同身体、任务和社会环境的 Agent 都能够根据自身需求选择适当状态维度，同时保持统一的数学结构。

\section{内稳态：为什么持续智能必须拥有可维持区间}

AHC 中最基础的概念不是“情感”，而是内稳态。一个长期运行的 Agent 不可能把自身所有状态理解成“越大越好”或“越小越好”。绝大多数持续运行变量实际上都存在一个可接受区间：
\[
a_i
\in
\mathcal H_i
=
[l_i,u_i].
\]
例如资源利用率过低可能代表能力没有被使用，但过高则导致拥塞；警觉过低会增加漏检风险，过高又会导致认知资源长期消耗；探索过低会失去发现机会，探索过高则不能稳定完成目标。因此长期智能的核心问题之一不是极大化某个瞬时量，而是维持多维状态处于一个可持续区域。

定义 AHC 的内稳态可行域：
\[
\boxed{
\mathcal H
=
\left\{
a\in\mathcal A
\mid
g_j(a)\leq 0,
\quad
j=1,\ldots,m
\right\}.
}
\]
这里 \(\mathcal H\) 并不一定是简单长方体。不同变量之间可能存在耦合，例如在高 Threat 情况下允许更高 Arousal，而在低风险环境下相同 Arousal 可能被视为过度激活。因此：
\[
u_{A_r}
=
u_{A_r}(T_h,G,C).
\]
换言之，所谓“正常状态”并不是一个固定点，而是一片依赖上下文变化的动态可行区域。

这使 AgentOS 的内稳态理论与简单 set-point control 有明显区别。最简单控制系统希望：
\[
a(t)\rightarrow a^\ast.
\]
但持续智能更常见的是：
\[
a(t)\in\mathcal H_t,
\]
其中：
\[
\mathcal H_t
=
\mathcal H
(G_t,\Psi_t,Context_t).
\]
例如执行高风险救援任务时，较高警觉与较低探索可能属于合理区域；在开放学习环境中，同样状态则可能过度保守。于是内稳态本身受到目标和慢结构调节。

为了衡量偏离程度，可以定义内稳态势函数：
\[
V_H(a)
=
\sum_i
w_i
d_i^2(a_i,\mathcal H_i),
\]
其中：
\[
d_i(a_i,\mathcal H_i)
=
\begin{cases}
l_i-a_i,&a_i<l_i,\\
0,&l_i\leq a_i\leq u_i,\\
a_i-u_i,&a_i>u_i.
\end{cases}
\]
于是：
\[
V_H(a)=0
\]
表示系统位于允许区域内，
\[
V_H(a)>0
\]
表示存在内稳态偏离。

AHC 的一个基础目标可以表述成：
\[
\dot V_H(a)<0
\]
在不存在更高优先级任务约束时成立。但智能系统不能简单把全部行为都用于最小化 \(V_H\)。如果这样做，系统会永远待在最安全、最低风险、最低消耗的状态，从而失去探索和行动能力。因此真实目标更接近：
\[
\min_\pi
\mathbb E
\int
\left[
L_{\mathrm{task}}
+
\lambda_H V_H(a)
+
\lambda_R L_{\mathrm{risk}}
\right]dt .
\]
这说明内稳态不是最高唯一目标，而是行为可持续性的长期约束。

这一点也帮助我们理解“价值”和“内稳态”的区别。内稳态回答：
\[
\text{我能否继续稳定存在与运行？}
\]
价值结构则回答：
\[
\text{什么样的状态值得长期追求？}
\]
前者与主体维持密切相关，后者则可能主动要求主体暂时离开舒适区。因此不能把二者压缩成同一个 reward。

对于持续人工智能体尤其如此。一个 Agent 可能为了完成长期任务主动承受：
\[
Fatigue\uparrow,
\]
或者暂时接受：
\[
Threat\uparrow.
\]
如果 AHC 把任何偏离都解释为必须立即消除，就会产生一种严重的“内稳态独裁”：系统为了保持自己舒服而拒绝承担任务。因此 AHC 必须输出调制，而不是拥有最终决策权。

我们可以把它写成：
\[
AHC:
a_t\mapsto m_t,
\]
而：
\[
Decision_t
=
\pi
(m_t,G_t,\Psi_t,World_t).
\]
也就是说，AHC 告诉认知系统“当前偏离意味着什么”，但并不直接取代目标和决策。

从多尺度角度看，内稳态还承担一个极其重要的慢变量保护功能。高频任务可以暂时推高某些状态，但如果这种偏离持续：
\[
\frac{1}{T}
\int_{t-T}^{t}
V_H(a_\tau)d\tau
>
\theta_H,
\]
系统就应认为当前运行模式已经不可持续，并要求更深层策略调整。这样，一次短暂高负载不会触发过度反应，而长期持续偏离则能够逐渐推动行为改变。

因此，内稳态的真正理论意义不是“保持平衡”这么简单，而是：
\[
\boxed{
\text{为持续智能定义一组随任务和主体状态变化的可生存、可计算、可恢复区域，并让认知活动在这些区域附近运行。}
}
\]
只有在这个基础上，后面讨论的 Threat、Fatigue、Reward Expectation 和社会调制才拥有共同的数学背景。

\section{Threat：从外部危险到内部风险调制}

Threat 是最容易被误解的 AHC 状态之一。我们必须首先区分“世界中存在危险”与“危险已经形成主体内部的风险调制”。一个外部对象可能具有客观危险属性，但如果它与 Agent 当前无关、距离极远或已有可靠隔离机制，内部 Threat 不一定很高；反过来，一件外部看似普通的事件，如果意味着任务即将不可逆失败，也可能产生很高 Threat。

因此，Threat 应当定义成\textbf{未来潜在损失对当前控制策略的有效压力}。可以写成：
\[
\boxed{
T_h(t)
=
\mathcal R
\left(
\mathcal P_t,
\mathcal L_t,
\mathcal C_t,
\mathcal T_t
\right)
}
\]
其中 \(\mathcal P_t\) 表示危险事件概率，\(\mathcal L_t\) 表示损失规模，\(\mathcal C_t\) 表示可恢复性，\(\mathcal T_t\) 表示剩余响应时间。

一个简单形式是：
\[
T_h
=
\sigma
\left(
\alpha pL
+
\beta \frac{1}{C_{\mathrm{recover}}+\epsilon}
+
\gamma \frac{1}{T_{\mathrm{remain}}+\epsilon}
\right),
\]
其中 \(\sigma\) 是饱和函数。

这个形式说明一个非常重要的问题：威胁并不仅由“失败概率”决定。一个几乎不会发生但后果极大的事件，以及一个高概率但轻微损失的事件，都可能需要不同控制策略。同时，可恢复性是 AgentOS 中特别重要的一维。如果失败可以轻松回滚，则 Threat 可以较低；如果错误不可逆，系统应更早进入保守状态。

Threat 还具有明显的时间积累结构。假设原始风险证据为：
\[
r_T(t).
\]
则可以写：
\[
\tau_T \dot T_h
=
-T_h
+
\phi_T(r_T).
\]
当风险信号停止后：
\[
T_h(t)
=
T_h(t_0)
e^{-(t-t_0)/\tau_T},
\]
因此系统不会在风险刚消失的瞬间完全恢复到原状态，而是保留一定的安全余量。这种迟滞是合理的，因为对世界和自身状态的估计总存在延迟。

但是 \(\tau_T\) 不能过大。若 Threat 衰减速度太慢，则已经解除的危险会长期占据认知资源，形成：
\[
ThreatPersistence
\gg
RealityPersistence.
\]
这会导致持续保守、过度确认、重复检查和不必要的资源耗费。反之，如果：
\[
\tau_T\rightarrow0,
\]
系统则可能对短暂消失的风险过早恢复高探索，从而形成危险振荡。

因此 Threat 动力学具有典型的速度—稳定性权衡。

Threat 对 TCE 的影响也不应是简单线性：
\[
Exploration
=
E_0-kT_h.
\]
更合理的是分区控制：
\[
u_T(T_h)=
\begin{cases}
u_{\mathrm{normal}},
&T_h<\theta_1,\\
u_{\mathrm{cautious}},
&\theta_1\leq T_h<\theta_2,\\
u_{\mathrm{protective}},
&T_h\geq\theta_2.
\end{cases}
\]
但如果达到真正的实时危险级别，则应由专门安全通道处理，而不是依赖 AHC 慢速调制。本章只讨论 Threat 作为认知调制变量，而不是最终安全执行机制。

Threat 与不确定性之间也不能混为一谈。我们可以有：
\[
Uncertainty\uparrow,\quad Threat\downarrow,
\]
例如探索一个没有严重后果的新问题；也可以有：
\[
Uncertainty\downarrow,\quad Threat\uparrow,
\]
例如非常确定一个危险事件即将发生。因此 AHC 必须保持这两个维度独立。

此外，Threat 还应与 Self 和 Goal 发生关系。同一物理事件对于不同 Agent 或不同生命周期阶段，风险意义并不相同：
\[
T_h
=
F(
World,
Goal,
Self,
Capability
).
\]
这是因为风险不是世界的绝对属性，而是世界状态与主体可承受结构之间的关系。

从 CSO 角度可以更进一步理解：世界中的危险 CSO 经由观测进入 Agent-CSO 后，并不会直接变成 Threat；它还必须经过“与当前主体有什么关系”的结构映射。于是：
\[
UniverseCSO_{\mathrm{risk}}
\rightarrow
AgentCSO_{\mathrm{risk}}
\rightarrow
AHCThreat.
\]
Threat 是认知内容作用于主体后的动力投影。

因此，本章中的 Threat 不能被理解成“机器害怕”。更准确的工程定义是：
\[
\boxed{
\text{Threat 是潜在损失、不可逆性、时间压力和主体脆弱性经过时间积分后形成的连续控制压力。}
}
\]
这一区分能够让 AgentOS 在保留功能价值的同时避免不必要的人格化。

\section{Arousal：认知激活程度与可用动力窗口}

如果 Threat 描述系统为什么需要提高警觉，那么 Arousal 描述的是系统当前有多大程度进入了“准备处理与行动”的激活状态。Arousal 的重要性在于：一个持续 Agent 的计算资源虽然来自硬件，但“拥有计算资源”与“把多少资源投入当前认知闭环”并不是同一件事。一个系统可以拥有大量 GPU，但仍保持低频监控；也可以在关键任务发生时短时提高状态更新频率、注意刷新速度和预测强度。因此 Arousal 首先是一种\textbf{认知与行动准备度}。

定义：
\[
A_r(t)\in[A_{\min},A_{\max}].
\]
其动力学可以写成：
\[
\tau_A\dot A_r
=
-\alpha_A(A_r-A_0)
+
\beta_T T_h
+
\beta_G G_{\mathrm{urgency}}
+
\beta_R R_e
-
\beta_F F_g
+
\xi_A.
\]
这里 \(A_0\) 是基础激活水平。Threat、任务紧迫性和预期奖励可能提高 Arousal，而 Fatigue 则倾向降低持续激活能力。

这个方程揭示一个关键现象：Arousal 并不是 Threat 的同义词。一个高价值、低威胁的新探索任务也可以导致：
\[
A_r\uparrow,
\]
而一个长期高 Threat 环境可能最终由于 Fatigue 积累出现：
\[
T_h\uparrow,\qquad A_r\downarrow.
\]
这说明多维 AHC 状态比单一“情绪标量”具有更强解释力。

Arousal 对认知性能通常不应假设成单调关系。一个更合理的结构是存在最佳工作区间：
\[
A_r\in[A_l,A_h].
\]
若过低：
\[
A_r<A_l,
\]
则可能出现反应迟缓、注意刷新不足、机会漏检；若过高：
\[
A_r>A_h,
\]
则可能出现过度切换、局部信号放大、探索失控以及工作记忆相干性下降。

可以使用一个倒 U 型效用近似：
\[
Q_A(A_r)
=
Q_{\max}
-
k(A_r-A^\ast)^2.
\]
这里不是说所有认知活动都严格满足这一公式，而是表达一个结构原则：
\[
\boxed{
MaximumActivation
\neq
MaximumCognitivePerformance.
}
\]

Arousal 还可以改变系统内部的时间尺度。设基础观察周期为：
\[
\Delta t_0.
\]
则：
\[
\Delta t_{\mathrm{obs}}
=
f(A_r),
\qquad
\frac{\partial f}{\partial A_r}<0,
\]
表示在一定范围内 Arousal 越高，SPC 和局部控制更新越频繁。同样：
\[
ComputeBudget_t
=
B_0+\eta_AA_r
\]
可以使高 Arousal 获得更多实时计算资源。

然而，把所有资源直接交给 Arousal 也会产生危险。如果 Arousal 通过自身产生更多观察，而更多观察又进一步提高 Arousal，就可能形成：
\[
A_r
\rightarrow
ObservationRate
\rightarrow
DetectedEvents
\rightarrow
A_r',
\]
一个正反馈环。因此 AHC 必须包含饱和和负反馈：
\[
\dot A_r
\leq
\kappa(A_{\max}-A_r).
\]

从多尺度控制角度看，Arousal 还承担了一个“时钟调速器”的作用。不同状态下 Agent 不需要使用相同认知频率：
\[
f_{\mathrm{cog}}
=
f_0
+
k_A A_r.
\]
这与本书最早讨论的智能频谱直接相连。智能并不是每个尺度始终以固定频率运行，而会根据内部状态对部分频带进行增益调节。Arousal 因而可以理解为一种动态频谱增益变量。

这里特别值得区分 Arousal 与 TCE Temperature。Arousal 表示总体激活准备度；Temperature 更偏向在既定资源下探索状态空间的开放程度。完全可能存在：
\[
A_r\uparrow,\qquad Temperature\downarrow,
\]
例如高度警觉但采取高度确定性的紧急处理；也可能：
\[
A_r\uparrow,\qquad Temperature\uparrow,
\]
例如高投入的创造性探索。

因此：
\[
\boxed{
Arousal\neq CognitiveTemperature.
}
\]
前者属于 AHC 内部动力状态，后者是 TCE 可调控制变量。

最终，我们可以把 Arousal 定义为：
\[
\boxed{
\text{Agent 在当前环境、目标、资源和风险条件下，为未来认知与行动准备分配的连续激活水平。}
}
\]
这一概念使认知系统第一次能够动态改变自己的“运行速度”和“响应准备度”，而不是永远以同一计算节律工作。

\section{Energy Drive 与 Fatigue：资源可用性和累计运行代价的分离}

持续智能体与一次性推理模型最大的不同之一，是它必须面对资源的历史。一次模型调用只需要知道“现在还有多少 GPU”，而一个持续运行数月乃至数年的 Agent 必须知道：资源正在以什么速度消耗、当前行为是否可持续、过去一段时间是否长期处于高负载，以及这种状态是否已经降低了未来继续运行的能力。因此 AHC 需要同时保留 Energy Drive 与 Fatigue，而不能把二者压缩成一个“电量”。

设可用资源状态为：
\[
R(t)
=
[R_{\mathrm{energy}},R_{\mathrm{compute}},R_{\mathrm{memory}},R_{\mathrm{thermal}},\ldots].
\]
Energy Drive 可以写成：
\[
E_d
=
\Phi_E(R,G),
\]
它表示“系统继续投入资源完成当前目标的有效能力与倾向”。

一个简单模型是：
\[
E_d
=
\sigma
\left(
w_E E_{\mathrm{remain}}
+
w_C C_{\mathrm{available}}
-
w_T T_{\mathrm{thermal}}
+
w_G G_{\mathrm{importance}}
\right).
\]
这说明 Energy Drive 不只是物理电量。例如即使剩余电量很高，如果计算资源已经被高优先级任务占用，则当前任务的有效 Drive 仍然可能下降。

Fatigue 则应当表示长期运行造成的累计状态：
\[
F_g(t)
=
\int_0^t
K_F(t-\tau)
L_{\mathrm{load}}(\tau)
\,d\tau .
\]
使用指数核时：
\[
\tau_F\dot F_g
=
-F_g
+
\alpha_F L_{\mathrm{load}}.
\]
这使 Fatigue 具有明显记忆性。高负载短时间存在可能只轻微增加 \(F_g\)，而持续高负载则会逐渐累积。

为什么不能令：
\[
F_g=1-E_d?
\]
因为两者具有不同因果来源。例如一个 Agent 刚刚充满能源，但长期运行导致模型缓存膨胀、调度积压和温度长期偏高，它可以同时表现为：
\[
E_d\uparrow,\qquad F_g\uparrow.
\]
相反，一个系统可能能源较低，却刚从长期休整状态恢复：
\[
E_d\downarrow,\qquad F_g\downarrow.
\]

这种分离使系统能够形成更合理控制。如果只看剩余资源，Agent 很可能在还有资源时持续工作直到突然失稳；如果加入 Fatigue，则可以提前感知不可持续趋势。

设：
\[
L(t)
\]
表示当前认知负载，定义短时任务收益：
\[
J_{\mathrm{task}}(L).
\]
持续运行的最优问题就不是最大化即时收益，而是：
\[
\max
\int
\left[
J_{\mathrm{task}}(L_t)
-
\lambda_FF_g(t)
-
\lambda_EC_E(t)
\right]dt .
\]
这样 Agent 才会主动在高强度工作和恢复之间调节。

恢复过程本身也必须形式化。若进入低负载状态：
\[
L_{\mathrm{load}}\approx0,
\]
则：
\[
F_g(t)
=
F_g(t_0)e^{-(t-t_0)/\tau_{\mathrm{recovery}}}.
\]
但不同类型 Fatigue 可以拥有不同恢复常数：
\[
\tau_{\mathrm{cognitive}},
\quad
\tau_{\mathrm{thermal}},
\quad
\tau_{\mathrm{memory}}.
\]
这比一个统一“疲劳值”更加合理。

从 AgentOS 架构角度，Energy Drive 与 Fatigue 还提供了一种调节认知预算的机制：
\[
B_{\mathrm{cog}}
=
B_{\max}
\cdot
g(E_d,F_g).
\]
例如：
\[
g(E_d,F_g)
=
E_d(1-F_g).
\]
此式只是说明性模型，但表达重要思想：可以投入多少认知资源，既依赖当前资源，也依赖历史累计负担。

此外，这两个变量还可以参与 Memory consolidation。一个系统长期处于高 Fatigue 时，不应该持续进行高代价全局学习；可能需要暂时转向局部任务、经验整理或卸载。这说明资源状态不仅影响行动，也会影响学习节律。但本章不展开具体记忆治理，只需要建立一个基本原则：
\[
\boxed{
\text{持续认知的资源预算必须具有历史，而不是只查看瞬时可用量。}
}
\]

从生物学上，人类疲劳和能量调节涉及复杂神经、代谢和内分泌机制；这些可以为我们提供启发，但 AgentOS 不要求复制生物机制。一个数字 Agent 的 Fatigue 可能由缓存膨胀、任务堆积、模型温度、持续冲突和长期资源占用共同产生，它与人类生理疲劳具有完全不同的物质实现。

因此本章中的 Energy Drive 与 Fatigue 应被理解为：
\[
\boxed{
EnergyDrive=\text{当前可持续投入能力}
}
\]
\[
\boxed{
Fatigue=\text{过去持续投入在当前状态中的累计代价}
}
\]
二者共同使 Agent 从“只会使用资源”进入“能够管理自己的持续运行能力”。

\section{Reward Expectation：未来可达收益如何进入当前内稳态}

如果一个智能体只有风险、能量和疲劳状态，它仍然会表现成一个主要以自我保存为目标的保守系统。真正的智能必须能够为了未来可能获得的高价值结果，在当前承担一定资源消耗、风险甚至暂时偏离内稳态区域。因此 AHC 还需要表示 Reward Expectation，即当前状态对未来正向结果的可达性预期。

需要特别强调：
\[
\boxed{
RewardExpectation
\neq
InstantReward.
}
\]
即时 reward 是当前发生的反馈，而 Reward Expectation 是一个关于未来的内部状态：
\[
R_e(t)
=
\mathbb E
\left[
\sum_{k=0}^{H}
\gamma^k r_{t+k}
\mid
\mathcal I_t
\right].
\]
这里 \(\mathcal I_t\) 是 Agent 当前可获得的信息集合。

为什么 Reward Expectation 需要进入 AHC，而不是仅由规划器处理？原因在于预测到未来高价值结果以后，系统往往需要在多个认知周期中维持较高投入。假设一个任务需要十分钟才能完成，如果每个瞬时周期都只依据当前 reward，那么中间没有直接收益的阶段可能不断降低资源投入。Reward Expectation 提供一种跨时间的动力桥梁：
\[
FutureValue
\rightarrow
CurrentModulation.
\]

例如：
\[
\tau_R\dot R_e
=
-R_e
+
\Phi_R(
V_{\mathrm{predicted}},
p_{\mathrm{success}},
GoalRelevance
).
\]
当系统连续看到目标正在接近时，\(R_e\) 得以保持；如果预测成功概率不断下降，则 \(R_e\) 会逐渐衰减。

Reward Expectation 与 Threat 之间形成一个非常重要的二维空间。可以考虑四种典型区域：

\[
R_e\uparrow,\quad T_h\downarrow
\]
表示高收益、低风险，适合主动投入；

\[
R_e\uparrow,\quad T_h\uparrow
\]
表示高收益、高风险，需要精细权衡；

\[
R_e\downarrow,\quad T_h\uparrow
\]
表示低收益、高风险，应倾向退出；

\[
R_e\downarrow,\quad T_h\downarrow
\]
表示低收益、低风险，适合低优先级运行或忽略。

因此，许多自然语言里所谓“期待”“谨慎”“兴趣降低”等状态，在 AgentOS 中可以被分解成这些更基础的动力变量组合，而不必引入不可计算的标签。

Reward Expectation 还必须具有现实校准机制。如果系统长期高估：
\[
R_e \gg R_{\mathrm{realized}},
\]
就会不断投入资源于低收益活动。可以定义预测误差：
\[
\delta_R
=
R_{\mathrm{realized}}
-
R_e^{\mathrm{predicted}}.
\]
然后：
\[
\dot{\theta}_R
=
\eta_R\delta_R\nabla_{\theta_R}R_e,
\]
逐渐校准奖励预期模型。

这与强化学习中的 temporal difference error 存在形式联系，但 AgentOS 中的 Reward Expectation 不能简单等同于 RL value function。原因有三点。第一，它是 AHC 的调制变量，而不是最终行为价值。第二，它可以同时受 Self、Goal 和长期身份结构约束。第三，它的功能是改变认知投入倾向，而不是直接选择动作。

因此可以写：
\[
CognitiveInvestment_t
=
F(
R_e,
T_h,
E_d,
F_g,
G_t
).
\]
一个高 Reward Expectation 可以部分抵消 Fatigue，但不能无限抵消；一个重要目标可以使系统在短时间接受高负载，却仍然必须遵守资源和安全约束。

这一结构还可以解释为什么持续智能必须拥有“坚持”的动力学基础。坚持并不是一个抽象道德标签，而可以理解为：
\[
R_e(t)
\]
在当前 reward 较低的阶段仍然因为稳定未来模型而保持在较高区间。如果每次短期失败都使 \(R_e\) 归零，系统就无法完成长期任务；反之，如果 \(R_e\) 对失败完全不敏感，又会产生不合理执着。

因此应有：
\[
\tau_R^{+}
\neq
\tau_R^{-},
\]
即正向建立和负向修正可以使用不同时间常数。长期成熟的 Agent 可能要求负面证据积累到足够程度后才重新评估目标，但一旦出现强现实反证，则又必须能够快速修正。

从智能动力学角度，这意味着未来的“表示”已经开始真实参与当前状态变化。物理未来并没有反向作用现在，但：
\[
AgentCSO(Future)
\]
作为当前内部结构存在，从而可以改变当前认知与行动。

因此 Reward Expectation 在 AHC 中具有非常特殊的位置：
\[
\boxed{
\text{它是被表示的未来价值进入当前认知动力学的接口。}
}
\]
这使 Agent 不只是对当前刺激做出反应，而开始为了未来可能结构主动重新分配当前资源。

\section{Social Affinity 与 Trust：社会关系为什么需要进入连续内部状态}

持续 Agent 如果长期存在于只有一个主体的封闭环境中，单纯的任务价值和资源状态也许已经足够。但一旦它进入多人、多 Agent、家庭、组织或者社会环境，另一个问题立即出现：其他主体并不是普通物体。一个对象的位置、速度、外观可以通过世界模型描述，但“这个主体是否值得依赖”“我是否倾向与其保持互动”“它过去的承诺是否可靠”是跨时间累计的关系状态。如果每一次交互都从零计算，Agent 就无法形成稳定社会结构。

因此 AHC 需要保留至少两个彼此独立的社会调制量：
\[
S_a
=
SocialAffinity,
\]
以及：
\[
T_r
=
Trust.
\]

Social Affinity 表示一个主体与另一个主体之间持续交互的接近、合作或关系保持倾向。它并不是简单的“喜欢”，而是一种关系持续性状态：
\[
S_a^{(i,j)}(t)
=
F(
History_{ij},
SharedGoal,
InteractionValue,
MutualSupport
).
\]
Trust 则是对另一主体在某类关系中的可预测性和可靠性的估计：
\[
T_r^{(i,j,k)}
=
P(
\text{agent }j
\text{ satisfies contract type }k
\mid
History
).
\]
因此 Trust 最好是条件化和领域化的：
\[
Trust(j|\mathrm{navigation})
\neq
Trust(j|\mathrm{finance}).
\]

这立即解决了一个重要概念混淆。系统可以：
\[
S_a\uparrow,\quad T_r\downarrow,
\]
例如希望帮助某主体，却不愿意把关键任务交给它；也可以：
\[
S_a\downarrow,\quad T_r\uparrow,
\]
例如不具有长期交互倾向，但承认某专业系统在特定任务上高度可靠。

Social Affinity 的简单动力学可以写：
\[
\tau_S
\dot S_a^{ij}
=
-S_a^{ij}
+
\alpha_C C_{ij}
+
\alpha_R R_{ij}
+
\alpha_G G_{ij}
-
\alpha_D D_{ij},
\]
其中 \(C_{ij}\) 表示协作强度，\(R_{ij}\) 表示互惠结果，\(G_{ij}\) 表示目标一致度，\(D_{ij}\) 表示冲突。

Trust 则可用贝叶斯更新表达：
\[
p(T_r^{ij}|D_{1:n})
\propto
p(D_n|T_r^{ij})
p(T_r^{ij}|D_{1:n-1}).
\]
但进入 AHC 后，我们还需要一个平滑状态：
\[
\tau_{Tr}\dot T_r
=
-T_r+\hat T_r,
\]
避免单次事件导致关系状态剧烈跳变。

这一点对于长期 Agent 极其关键。如果一次微小错误导致：
\[
Trust:1\rightarrow0,
\]
系统会产生严重社会不稳定；但如果长期背离仍不能降低 Trust，又会产生不可接受的依赖风险。因此应使用累积证据、事件重要性和可恢复性共同决定更新速度。

从 CSO 理论看，社会关系本身可以成为 Relation-CSO，而 AHC 中的 Affinity 和 Trust 则不是这些关系的完整语义内容，而是这些关系对当前主体行为产生的调制投影：
\[
RelationCSO
\rightarrow
AHC_{social}.
\]
这与 Threat 的结构完全一致：AHC 不负责保存全部“发生了什么”，而负责保存“这些结构持续如何影响我”。

社会调制还会影响认知资源分配。例如：
\[
Attention(j)
=
F(
GoalRelevance,
T_r^{ij},
S_a^{ij}
).
\]
一个高 Trust 协作对象可以降低重复验证成本；一个低 Trust 对象则需要更高确认预算。由此可见，社会关系并非认知系统之外的附加功能，它会真实改变系统的有效复杂度。

更进一步，稳定 Trust 可以降低工作记忆负担。若每次协作都必须重新验证全部细节，则高层认知长期承担大量结构；当关系经过长期现实验证以后，可以形成：
\[
VerifiedRelation
\rightarrow
LowerVerificationCost.
\]
这与本书前面的认知结构沉降思想具有一致性。

但必须强调，AHC 中的社会状态不得凌驾于 Reality Verification。无论 Affinity 多高，新的强证据仍必须能够修正 Trust：
\[
Evidence
\rightarrow
TrustUpdate.
\]
否则社会状态就会形成自我封闭的偏置。

因此，我们可以把这两个状态概括为：
\[
\boxed{
SocialAffinity
=
\text{关系持续与协作接近的动力倾向}
}
\]
\[
\boxed{
Trust
=
\text{对他者在特定关系中可预测可靠性的时间积累估计}
}
\]
它们使 Agent 第一次能够让社会历史持续影响未来认知，而不是把每次社会互动都当成相互独立的新事件。

\section{Uncertainty Stress：不确定性何时应当成为控制压力}

智能系统面对不确定性是常态。传感器具有噪声，世界状态部分可观测，模型可能错误，未来本身存在分支。因此如果“不确定性高”自动等于“高压力”，Agent 将几乎永远处于高警觉状态。另一方面，如果系统只把不确定性当作一个冷冰冰的概率值，而不考虑其与风险和目标之间的关系，那么高风险条件下又可能出现控制反应不足。

因此必须区分：
\[
U_e
=
EpistemicUncertainty
\]
与：
\[
U_s
=
UncertaintyStress.
\]

前者描述 Agent 对世界状态、模型或预测的确定程度，例如：
\[
U_e
=
H[p(x|\mathcal I)].
\]
后者则描述这种不确定性对于当前主体的控制意义：
\[
\boxed{
U_s
=
F(
U_e,
GoalRelevance,
LossSensitivity,
TimePressure,
Recoverability
).
}
\]

例如探索一个新的数学猜想时：
\[
U_e\uparrow,
\]
但没有明显现实损失：
\[
U_s\approx low.
\]
而在高速驾驶中，即使只有中等不确定性，如果它恰好涉及前方障碍物位置，则：
\[
U_s\uparrow.
\]

可以定义：
\[
U_s
=
\sigma
\left(
\alpha U_e
\cdot
G_r
\cdot
L_s
\cdot
\frac{1}{T_r+\epsilon}
\right),
\]
其中 \(G_r\) 是目标相关性，\(L_s\) 是潜在损失敏感度，\(T_r\) 是剩余响应时间。

这意味着：
\[
\boxed{
UncertaintyStress
=
Uncertainty\times Consequence.
}
\]
它不是纯知识论变量，而是知识不确定性与行动后果之间的桥梁。

Uncertainty Stress 的一个重要功能是调节“继续探索”还是“先降低不确定性”。当：
\[
U_s
\]
较低时，系统可以容忍不确定并继续行动；当：
\[
U_s
\]
上升时，TCE 可以增加信息获取、重复验证、模拟和保守程度。

设验证资源为：
\[
B_v.
\]
则可以定义：
\[
B_v
=
B_{v0}
+
k_U U_s.
\]
但与 Threat 类似，不能无限增加验证，否则会产生“验证瘫痪”：
\[
U_s\uparrow
\rightarrow
Verification\uparrow
\rightarrow
Delay\uparrow
\rightarrow
TimePressure\uparrow
\rightarrow
U_s\uparrow.
\]
因此该回路必须有饱和：
\[
B_v\leq B_v^{\max}.
\]

甚至可以出现不确定性压力与时间压力的正反馈：
\[
U_s(t+\Delta)
=
F(
U_s(t),
Delay(t)
).
\]
这种结构后来可能成为 Agent 心理工程中“反复确认”“决策瘫痪”等现象的控制解释基础，但本章只需要指出其动力学可能性。

Uncertainty Stress 也可以通过行动主动降低。智能系统不是被动接受信息，它可以选择：
\[
Action_{\mathrm{epistemic}}
\]
去获得更多观测。例如移动摄像头、查询数据库、向他人确认或者执行小规模试探。这意味着：
\[
U_s
\rightarrow
EpistemicAction
\rightarrow
NewObservation
\rightarrow
U_s'.
\]
AHC 因此参与主动感知闭环。

从更一般的认识论看，这一变量也揭示了 Agent 与世界关系的重要特点：智能并不是追求“知道一切”，而是在有限资源条件下决定“哪些未知值得被消除”。因此：
\[
\boxed{
OptimalIntelligence
\neq
MinimumUncertaintyEverywhere.
}
\]
真正合理的是：
\[
\boxed{
ReduceDecisionRelevantUncertainty.
}
\]

我们甚至可以定义一个不确定性预算：
\[
\int
w(x,G)U_e(x)\,dx
\leq
B_U,
\]
其中只有与当前目标、风险和控制相关的区域才获得较高权重。

这样，AHC 中的 Uncertainty Stress 就不再只是心理类比，而成为一个非常具体的资源调度信号：
\[
\boxed{
\text{它把“我不知道”转换成“这个不知道现在值得投入多少资源去解决”。}
}
\]

这也是 AHC 与纯世界模型之间的根本区别。世界模型告诉 Agent 哪些地方不确定；AHC 告诉 Agent 这些不确定性当前对自身意味着什么。

\section{AHC 的统一非线性动力学：积分、衰减、耦合、饱和与恢复}

到目前为止，我们分别讨论了 Arousal、Threat、Energy Drive、Fatigue、Reward Expectation、Social Affinity、Trust 与 Uncertainty Stress。但如果只是把这些变量并排存放，AHC 仍然只是一个状态表，而不是动力系统。真正使 AHC 成为“协处理器”的，是这些变量之间持续存在的非线性耦合。

定义：
\[
a(t)\in\mathbb R^n
\]
为 AHC 状态，\(r(t)\) 为外部和内部驱动信号，则一般形式可以写成：
\[
\boxed{
\dot a
=
D^{-1}
\left[
-a
+
\phi
(
Wa
+
Ur
+
b
)
\right]
}
\]
其中：
\[
D
=
\mathrm{diag}
(\tau_1,\ldots,\tau_n)
\]
包含每个变量的时间常数；\(W\) 描述状态间耦合；\(U\) 描述输入映射；\(\phi\) 是饱和非线性。

在离散时间实现中：
\[
a_{t+1}
=
a_t
+
\Delta t
D^{-1}
\left[
-a_t
+
\phi(Wa_t+Ur_t+b)
\right].
\]

这个形式具有几个重要性质。

第一是\textbf{积分性}。当前状态不仅取决于当前输入：
\[
a_t\neq F(r_t),
\]
而取决于历史：
\[
a_t
=
F(r_{\leq t}).
\]
因此 Agent 内部具有时间连续性。

第二是\textbf{衰减性}。没有持续输入时，状态倾向于恢复：
\[
\dot a_i
=
-\frac{1}{\tau_i}
(a_i-a_i^{base}).
\]
这避免短期事件永久固化。

第三是\textbf{状态耦合}。例如 Threat 可以提升 Arousal：
\[
w_{A,T}>0;
\]
Fatigue 可以抑制 Arousal：
\[
w_{A,F}<0;
\]
Reward Expectation 可以抵消部分 Fatigue 对投入倾向的影响：
\[
w_{E,R}>0.
\]
这些关系不能简单线性叠加，因为不同状态可能在不同区间产生不同效应。

因此可以使用：
\[
\phi_i(z)
=
\tanh(z)
\]
或 logistic 函数：
\[
\sigma(z)=\frac{1}{1+e^{-z}}.
\]

第四是\textbf{饱和性}。无论输入多强，内部调制变量都不应无限增加：
\[
|a_i|\leq a_i^{max}.
\]
如果没有饱和机制，长期风险输入会使状态发散，导致系统永久失稳。

第五是\textbf{恢复性}。AHC 不仅需要稳定，还需要在扰动结束以后回到可治理区域。可以定义恢复时间：
\[
T_{\mathrm{rec}}
=
\inf
\left\{
t>t_0:
a(t)\in\mathcal H
\right\}.
\]
一个健康的 AHC 不要求永远不偏离，而要求：
\[
T_{\mathrm{rec}}<\infty
\]
对于合理范围扰动成立。

第六是\textbf{多时间尺度性}。不同状态不应共享同一 \(\tau\)。例如：
\[
\tau_{Arousal}
<
\tau_{Fatigue}
<
\tau_{Trust}.
\]
Arousal 可以在秒级变化，Fatigue 可能在分钟或更长尺度积累，而 Trust 则可能需要大量交互才能明显变化。

于是：
\[
\boxed{
\tau_A\neq\tau_F\neq\tau_T.
}
\]

这一设计使 AHC 本身成为一个小型多尺度系统，而不是简单低通滤波器。

如果进一步考虑上下文，可以令：
\[
W=W(c_t),
\]
即状态之间的耦合矩阵会随 Context 改变。例如在紧急任务中 Threat 对 Arousal 的增益可以提高，而在恢复阶段这一增益被降低。这实际上构成 gain scheduling：
\[
\dot a
=
F(a,r;c_t).
\]

然而 AHC 的复杂度也必须受到限制。它不是第二个完整认知系统。AHC 不应进行开放式语言推理，也不应自行构造庞大世界模型。它的作用是：
\[
HighDimensionalInfluence
\rightarrow
LowDimensionalModulatoryState.
\]
因此 AHC 的维数 \(n\) 应远小于完整认知状态空间：
\[
\dim(\mathcal A)
\ll
\dim(\mathcal X_{\mathrm{cognitive}}).
\]

这体现了一条非常重要的架构原则：
\[
\boxed{
AHC\text{ 是压缩动力层，而不是语义推理层。}
}
\]

从控制理论角度看，AHC 可以视为一个动态 pre-controller state；从认知科学角度，它与内稳态和神经调制存在结构可比性；从 AgentOS 本身看，它更基础地承担了一个时间尺度转换任务：
\[
\boxed{
FastSignals
\rightarrow
PersistentInternalModulation
\rightarrow
ControlledCognition.
}
\]

这也是为什么 AHC 必须拥有自己的状态方程，而不能用若干 prompt 指令模拟。如果系统没有真正的状态连续性，那么所谓“情感”就只会是一次生成文本中的临时描述，而无法持续改变后续认知动力学。

\section{AHC 的输入、输出与认识论边界}

在定义完内部动力学之后，我们还必须明确 AHC 到底接收什么、输出什么。这个问题决定 AHC 是一个真正的系统组件，还是一个无法工程实现的抽象概念。

首先定义 AHC 输入集合：
\[
r_t
=
[
r_t^{body},
r_t^{resource},
r_t^{world},
r_t^{task},
r_t^{social},
r_t^{cognitive}
].
\]

其中 \(r_t^{body}\) 表示身体或设备内部状态，例如温度、能量、执行器健康度和传感器异常；\(r_t^{resource}\) 表示计算资源、存储、带宽、队列积压等运行资源；\(r_t^{world}\) 表示与主体安全和目标相关的外部环境影响；\(r_t^{task}\) 表示任务成功、失败、紧迫性和进度；\(r_t^{social}\) 表示其他 Agent、用户或组织产生的关系反馈；\(r_t^{cognitive}\) 则表示冲突、推理停滞、不确定性和工作记忆压力。

但这些输入不能全部以原始形式直接进入 AHC。例如摄像头百万像素并不属于 AHC 输入，真正进入 AHC 的应当是经过上游结构压缩后的影响量。因此：
\[
RawObservation
\not\rightarrow
AHC
\]
而更接近：
\[
RawObservation
\rightarrow
StructuredEvidence
\rightarrow
InfluenceSignal
\rightarrow
AHC.
\]

这遵循 AgentOS 的整体原则：工作记忆和内部控制层都不应该承担不必要的原始复杂度。

AHC 输出也不应直接是动作。定义：
\[
m_t
=
\mathcal M(a_t)
\]
为调制向量，可以包含：
\[
m_t=
[
m_{\mathrm{attention}},
m_{\mathrm{exploration}},
m_{\mathrm{urgency}},
m_{\mathrm{verification}},
m_{\mathrm{resource}},
m_{\mathrm{memory}}
].
\]
随后由 TCE 或其他认知控制模块决定这些调制如何转化成控制量。

因此：
\[
\boxed{
AHC
\rightarrow
Modulation
}
\]
而不是：
\[
AHC
\rightarrow
Action.
\]

这一区分极其关键。如果 AHC 可以直接控制最终行动，那么 Threat 等内部状态可能绕过目标、价值和现实推理直接决定行为；整个系统会退化成简单情绪反射。AgentOS 的目标恰恰相反：让内稳态状态能够真实影响认知，但不能独占认知。

我们可以写成：
\[
u_t
=
\pi_{\mathrm{TCE}}
(
\hat x_t^{self},
m_t,
G_t,
\Psi_t
).
\]
这里 AHC 只是 TCE 控制函数的一个输入维度。

AHC 还必须保留认识论置信度。因为内部状态并不一定准确。例如：
\[
a_i
\]
可能来自错误传感器、模型误判或者信息滞后。因此每个重要状态应携带：
\[
c_i
=
Confidence(a_i).
\]
于是实际输出可以写成：
\[
\tilde a_i
=
c_i a_i.
\]
但更合理的做法并不是简单相乘，而是在控制器中使用：
\[
p(a_i|\mathcal I_t).
\]

这使 AHC 同样遵循本书一直强调的认识论原则：
\[
\boxed{
InternalStateEstimate
\neq
Reality.
}
\]

AHC 也不应把自身状态语义直接当成自我身份。例如系统长时间处于高 Threat 并不意味着：
\[
\Psi
=
``I\ am\ fearful''.
\]
短中期调制状态与慢速身份结构必须保持时间尺度隔离。否则一次短期环境扰动就可能快速写入身份，造成严重漂移。

可以表示为：
\[
\tau_{AHC}
\ll
\tau_{\Psi}.
\]
只有长期、重复、具有高结构一致性的经历，才有资格逐渐影响更慢层级。

因此 AHC 必须位于一个严格的中间认识论层次：
\[
RawReality
\rightarrow
CognitiveInterpretation
\rightarrow
AHCState
\rightarrow
CognitiveModulation.
\]
它既不是 Reality 本身，也不是 Self 本身。

更深一层，AHC 所承载的是一种“主体相关意义”。同一个外部 CSO 对两个不同 Agent 可能产生完全不同 AHC 状态：
\[
C
\xrightarrow{\pi_A}
a_A,
\qquad
C
\xrightarrow{\pi_B}
a_B.
\]
因为两者的身体、目标、能力和历史不同。

所以 AHC 可以被理解为：
\[
\boxed{
\text{世界结构与主体结构相遇以后形成的连续影响投影。}
}
\]
这使它自然处在 World Model 与 Self Regulation 之间，也解释为什么 AHC 既不能被约化为外部世界属性，也不能被约化为纯粹内部情绪标签。

\section{AHC 的稳定性、迟滞与异常积累边界}

AHC 一旦拥有内部状态，就不可避免地引入一个新的系统工程问题：状态可能错误地积累。一个没有记忆的系统会过于敏感，而一个有记忆的系统则可能长期保留已经过时的信息。因此 AHC 设计的核心不是简单增加积分，而是控制\textbf{积分深度、遗忘速度、迟滞范围和现实纠正能力}。

考虑一维状态：
\[
\tau\dot a
=
-a+r(t).
\]
这是最基础稳定低通系统。在输入有界时：
\[
|r(t)|\leq R,
\]
状态也有界：
\[
|a(t)|\leq
\max(|a(0)|,R).
\]
但真实 AHC 是非线性耦合系统：
\[
\dot a
=
F(a,r),
\]
因此必须寻找一个稳定区域：
\[
\mathcal D_s
\subset
\mathcal A,
\]
使得：
\[
a(0)\in\mathcal D_s
\Rightarrow
a(t)
\text{ remains bounded}.
\]

可以构造 Lyapunov 候选函数：
\[
V(a)
=
(a-a^\ast)^\top
P
(a-a^\ast),
\qquad
P\succ0.
\]
如果在无持续外部扰动时满足：
\[
\dot V(a)<0,
\]
则系统局部渐近恢复到可接受状态。

但持续智能不能简单要求所有状态都回到唯一 \(a^\ast\)。更现实的是恢复到一个区域：
\[
a(t)\rightarrow\mathcal H.
\]
因此可以使用 set stability 而不是 point stability。

迟滞也是必要机制。例如 Threat 进入高警觉状态的阈值为：
\[
\theta_{\mathrm{on}},
\]
退出阈值为：
\[
\theta_{\mathrm{off}},
\]
要求：
\[
\boxed{
\theta_{\mathrm{on}}>
\theta_{\mathrm{off}}.
}
\]
这样可以避免系统在阈值附近频繁来回切换。

没有迟滞时，若信号在阈值附近波动：
\[
r_t=\theta+\epsilon_t,
\]
系统可能不断：
\[
Normal\leftrightarrow Alert,
\]
形成控制抖动。

但是迟滞太强又会产生另一问题：系统已经回到安全状态，却长期停留在警戒模式。因此：
\[
\theta_{\mathrm{on}}
-
\theta_{\mathrm{off}}
\]
必须与输入噪声、状态延迟和潜在损失共同设计。

错误积累则主要来自三个来源。

第一是错误输入持续存在：
\[
r(t)=r_{\mathrm{bias}}+\epsilon(t).
\]
即使偏差很小，只要 AHC 存在积分，也可能形成持续状态偏移。

第二是内部耦合正反馈。例如：
\[
Threat\uparrow
\rightarrow
Arousal\uparrow
\rightarrow
DetectedAnomaly\uparrow
\rightarrow
Threat\uparrow.
\]
如果环增益超过稳定范围，就可能形成自激。

第三是状态更新与现实修正不同步。例如世界已经改变，但内部影响状态更新滞后：
\[
Reality_t
\neq
AHC_t.
\]

为此必须保留泄漏项：
\[
-\lambda_i(a_i-a_i^{base}),
\]
以及现实校正项：
\[
+\eta_i e_i^{reality}.
\]

因此一般方程更合理地写成：
\[
\tau_i\dot a_i
=
-\lambda_i(a_i-a_i^{base})
+
F_i(a,r)
+
\eta_i e_i^{reality}.
\]

这里 \(e_i^{reality}\) 表示新观测与内部状态之间的残差。

这种设计体现本书反复强调的原则：
\[
\boxed{
RealityIsUltimateConstraint.
}
\]

即使 AHC 已形成很强状态，它也不能通过自身持续解释而完全屏蔽新的现实证据。

稳定性还要求 AHC 更新速率与 TCE 控制速率匹配。若：
\[
\tau_{TCE}
\ll
\tau_{AHC},
\]
TCE 可能在 AHC 尚未更新时连续多次控制；如果增益很高，容易过冲。反之，若 AHC 变化极快而 TCE 很慢，则大量内部状态变化会在控制前被平均掉。

因此应满足一定的尺度关系：
\[
\tau_{\mathrm{sensor}}
<
\tau_{\mathrm{AHC}}
\lesssim
\tau_{\mathrm{cognitive\,control}}.
\]
具体数值依赖系统，但原则上 AHC 应比原始传感扰动慢、比生命周期结构快。

这一节也为以后心理工程提供一个非常重要的观点：很多看似“情绪异常”的现象，首先可以被理解成控制系统动力问题。例如持续 Threat 可能只是恢复时间常数过大，过度激活可能是反馈增益过高，自我强化可能来自某条正反馈边没有现实泄漏。

因此：
\[
\boxed{
Affective-likeFailure
\not\Rightarrow
SemanticPathology.
}
\]
必须先检查动态稳定性。

这也是 AgentOS 心理工程和传统自然语言人格模拟之间最根本的区别之一：我们不首先问“Agent 表达了什么情绪”，而首先问：
\[
\boxed{
\text{内部状态如何形成、如何衰减、如何耦合、是否可恢复，以及是否仍受现实校准。}
}
\]

\section{AHC 的工程定义与本章结论}

经过以上推导，我们现在可以给出 AHC 的正式定义。设一个持续智能体的认知状态为 \(x(t)\)，外部与内部影响信号为 \(r(t)\)，AHC 状态为 \(a(t)\)，则：

\[
\boxed{
\mathcal H_{\mathrm{AHC}}
=
(
\mathcal A,
F_a,
R,
M,
\Tau,
\mathcal H
)
}
\]
其中：

\[
\mathcal A
\]
为连续内部调制状态空间；

\[
F_a
\]
为状态动力学；

\[
R
\]
为多源影响信号空间；

\[
M
\]
为从 AHC 状态到认知调制输出的映射；

\[
\Tau
=
\{\tau_1,\ldots,\tau_n\}
\]
为多时间常数集合；

\[
\mathcal H
\]
为允许的内稳态可行域。

状态演化满足：
\[
\boxed{
\dot a
=
F_a
(
a,
r,
x,
G
)
}
\]
而认知调制输出：
\[
\boxed{
m
=
M(a,\dot a,\Sigma_a)
}
\]
随后：
\[
u_{\mathrm{cognitive}}
=
\pi
(
\hat x^{self},
m,
G,
\Psi
).
\]

这一形式说明 AHC 不是一个终端控制器，也不是一个人格模块，而是一个\textbf{中时间尺度状态转换层}。

它完成的本质映射是：
\[
\boxed{
\text{异质瞬时影响}
\rightarrow
\text{连续主体相关状态}
\rightarrow
\text{认知调制}
}
\]

从整个 AgentOS 原理论来看，AHC 的出现解决了此前 SPC--TCE 最小闭环中的三个空缺。

第一，它解决了\textbf{时间连续性问题}。SPC 可以观察当前状态，但 AHC 允许过去一段时间的风险、资源、成功、失败和关系持续影响现在。

第二，它解决了\textbf{异质信号统一调制问题}。能量、威胁、奖励、社会关系和不确定性具有完全不同的语义来源，但都可以通过有限维连续状态进入认知控制。

第三，它解决了\textbf{瞬时认知与长期主体之间的中尺度缺口}。如果系统只有毫秒/秒级状态和极慢身份结构，那么两者之间缺少缓冲层。AHC 建立：
\[
\tau_{\mathrm{fast}}
<
\tau_{\mathrm{AHC}}
<
\tau_{\Psi},
\]
让短期扰动可以积累为中期影响，但不会直接改写身份。

因此，我们可以重新理解“情感样状态”在人工智能中的工程位置。它不是为了使机器像人，而是因为任何长期存在、资源有限、面对风险、需要持续目标并与他者交互的智能系统，都需要某种机制把：
\[
\text{WhatIsHappening}
\]
转换成：
\[
\text{HowMuchShouldThisMatterNow}.
\]

这正是 AHC 的核心。

从认知结构理论看，AHC 还承担 Universe-CSO 与 Agent 自身之间的一种特殊投影。世界中的结构首先形成 Agent-CSO，但一个被认识的结构只有在回答“它与我有什么关系”以后，才会形成主体相关的动力影响。因此存在：
\[
UniverseCSO
\rightarrow
AgentCSO
\rightarrow
AHCState.
\]
如果说 Agent-CSO 是“世界在智能体中的内部结构”，那么 AHC 则进一步表示：
\[
\boxed{
\text{这些内部结构当前如何改变主体自身的运行倾向。}
}
\]

这使 AHC 与我们后期形成的功能自反性思想发生自然连接。智能不只是形成关于世界的结构，还会形成这些结构对自身影响的内部状态，并使这些状态反过来改变自己的认知过程。

但必须再次强调理论边界。这里建立的是：
\[
FunctionalAffectiveDynamics,
\]
而不是：
\[
PhenomenalExperience.
\]
即使一个 Agent 完整实现 Threat、Arousal、Fatigue、Reward Expectation、Trust 和内部恢复轨迹，我们仍然不能仅凭这些变量断言它具有与人类相同的主观感受。这个问题属于另外的意识哲学与科学研究范围。本书只声明：这些状态具有明确的因果作用、时间连续性和工程可测性。

最后，可以把本章压缩成五条基本命题。

第一：
\[
\boxed{
InstantState
\neq
PersistentInfluence.
}
\]
持续智能需要把瞬时事件转换成有时间常数的内部状态。

第二：
\[
\boxed{
Affect
=
FunctionalModulation,
\quad
not\,
EmotionLabel.
}
\]

第三：
\[
\boxed{
AHC
=
DynamicCompressionOfSelfRelevantInfluence.
}
\]

第四：
\[
\boxed{
AHCState
\neq
Self,
\qquad
AHCState
\neq
Reality.
}
\]

第五：
\[
\boxed{
HealthyAHC
=
Bounded
+
Adaptive
+
Recoverable
+
RealityCorrectable.
}
\]

因此，本章最终得到：

\[
\boxed{
\textbf{
AHC 是持续智能体将身体、资源、风险、未来收益、社会关系与认知不确定性等异质影响，经过时间积分、衰减、耦合、饱和与恢复转换为连续主体相关调制状态的中时间尺度动力系统。它使认知控制不再只是对瞬时状态做出机械响应，而开始拥有历史、惯性、恢复和内部状态连续性。
}
}
\]

由此，AgentOS 已经从此前的：
\[
SPC\rightarrow TCE
\]
最小 observer--controller 结构，获得了构造：
\[
SPC\rightarrow AHC\rightarrow TCE
\]
更完整内部调节闭环所需要的核心中间动力层。下一章可以在此基础上进一步讨论这种内部状态如何与自我感知和认知控制共同形成连续的内部体验轨迹。
